\begin{exercisechunk}
\item \textbf{Improving the worst context.}\suppmark
\label[exercise]{ex:l20-worst-context}
\exerciseprofile{3}{2}{3}{\exproof\quad\excalculation\quad\exconstruction\quad\exinterpretation}{%
Choose a direction that improves every tied worst context and obtain a convergence rate.}
Consider the finite binary population from \cref{sec:context-weighting},
with signed features $z_j$, unit separator $w$, and margin $\gamma>0$.
Put $p_j(\theta)=\sigma(z_j^\top\theta)$,
$Q(\theta)=\min_jp_j(\theta)$, and
$\cA(\theta)=\{j:p_j(\theta)=Q(\theta)\}$.
Let $C(\theta)$ be the convex hull of
$\{\nabla p_j(\theta):j\in\cA(\theta)\}$, and let $g(\theta)$ be its
minimum-norm point. Consider an absolutely continuous solution of
$\dot\theta=g(\theta)$ almost everywhere; existence may be assumed.
You may use the finite-minimum chain rule: along this flow,
$\frac{d}{dt}Q(\theta(t))=\min_{j\in\cA(\theta(t))}
\nabla p_j(\theta(t))^\top g(\theta(t))$ almost everywhere.
\begin{enumerate}[(a)]
\item Prove that $u^\top g\ge\|g\|_2^2$ for every $u\in C(\theta)$,
and that equality holds for at least one active gradient.
\item Writing $q(t)=Q(\theta(t))$, show that
$\dot q(t)=\|g(\theta(t))\|_2^2$ almost everywhere and
$\gamma q(t)(1-q(t))\le\|g(\theta(t))\|_2$.
\item Prove the rate in \cref{prop:worst-context}.
Give two equally likely, linearly separable contexts and a parameter
at which an arbitrary choice of one active gradient decreases the other
active probability to first order. Explain why the minimum-norm choice
avoids that problem.
\end{enumerate}
\begin{hints}
For (a), compare $\|g+\lambda(u-g)\|_2^2$ with $\|g\|_2^2$ for
$0\le\lambda\le1$. For (c), use the two signed features from the
lecture's PG failure example at zero initialization.
\end{hints}
\end{exercisechunk}
